% 本节给出单位周期立方孔附近的受限族结果，不是原 C_n 的全局上界。
\subsection{周期立方孔附近的小微扰不能提高二阶系数}
\label{sec:periodic-local-lower}

固定 $n\ge2$，记 $u=t^{1/n}$、$C_\delta=[-\delta,\delta]^n$，并令
\[
 D_\delta=\bigcup_{k\in\mathbb Z^n}(k+C_\delta),\qquad
 P=[-1/2,1/2)^n.
\]
取 $\mathbb Z^n$-周期缺陷集 $D$，它由 $D_\delta$ 删除或增加可测部分得到。以
\[
 T_0=(D_\delta\setminus D)\cap P,\quad
 S_0=(D\setminus D_\delta)\cap P,
 \qquad a=|T_0|,\quad p=|S_0|
\]
表示每胞删除和增加的部分，并记其周期延拓为
\[
 T=\bigcup_{k\in\mathbb Z^n}(k+T_0),\qquad
 S=\bigcup_{k\in\mathbb Z^n}(k+S_0).
\]
以下 $a,p$ 始终是每胞体积。令
\[
 \begin{gathered}
 E=\mathbb R^n\setminus D,\qquad H=\{M_{c,n}\ind_E>1-t\},\\
 q=|D\cap P|=(2\delta)^n-a+p,\qquad G=|H\cap D\cap P|,\\
 \mathcal P_t(D)=1+\frac{G}{1-q}.
 \end{gathered}
\]

\begin{theorem}[周期孔局部二阶排除]\label{thm:periodic-local-lower}
若 $\delta/u\to1$ 且对某个固定 $K<\infty$ 有 $a+p\le Ktu$，则
\[
 \limsup_{t\downarrow0}\frac{\mathcal P_t(D)-1-2^nt}{tu}
 \le -n2^{n+1}.
\]
该结论针对真实中心立方体极大算子的周期 halo，但只适用于单位周期立方孔附近、每胞对称差为 $O(tu)$ 的局部族，不能作为一般 $C_n$ 的上界。原立方孔 $a=p=0$、$\delta=u/(1+2u)$ 达到此二阶系数。
\end{theorem}

证明先控制可能产生新 halo 缺陷点的三类半径。置 $m_0=(2\delta)^n$，故 $m_0\sim2^nt$，$q/m_0\to1$。卫星缺陷的总量为每胞 $p$，可在最终捕获量中全额计入。

\begin{lemma}[内禀小半径估计]\label{lem:intrinsic-small-radius}
设 $\nu$ 是局部有限 Radon 测度，$A$ 可测且 $0<\nu(A)<\infty$。定义
\[
 M_{c,\nu}f(x)=\sup_{r>0,\ 0<\nu(Q(x,r))<\infty}
 \frac{1}{\nu(Q(x,r))}\int_{Q(x,r)}|f|\,d\nu.
\]
若对所有辅助 Borel 测度，中心方体极大算子的 weak $(1,1)$ 常数均不超过 $2^n$，则 $0<\gamma<1$ 时
\[
 \nu(\{M_{c,\nu}\ind_A>\gamma\}\setminus A)
 \le 2^n\frac{1-\gamma}{\gamma}\nu(A).
\]
\end{lemma}
\begin{proof}
令 $\beta=(1-\gamma)/\gamma$，取上述差集内有限测度紧子集 $K$，并定义
$d\sigma_c=(\ind_A+c\ind_K)d\nu$、$f=\ind_A$。在 $K$ 上的任一原见证 $Q$ 满足 $\nu(Q\setminus A)<\beta\nu(Q\cap A)$，从而
\[
 \frac{\int_Q f\,d\sigma_c}{\sigma_c(Q)}
 =\frac{\nu(Q\cap A)}{\nu(Q\cap A)+c\nu(Q\cap K)}
 >\frac1{1+c\beta}.
\]
对 $\sigma_c$ 使用假设的 weak 估计，得到
$c\nu(K)\le2^n(1+c\beta)\nu(A)$。除以 $c$，令 $c\to\infty$，再以内正则性逼近整个差集即得结论。
\end{proof}

\begin{proof}[定理~\ref{thm:periodic-local-lower} 的证明]
对 $x\in C_\delta$ 及 $r\le2\delta$，取 $d\nu=\ind_{D_\delta}dx$。每个坐标中 $Q(x,r)$ 与本孔的交长至少为 $r$，所以
\[
 \nu(Q(x,r))\ge r^n=2^{-n}|Q(x,r)|.
\]
若此方体见证 $x\in H$，则 $|D\cap Q|<t|Q|$，且
$D=(D_\delta\setminus T)\cup S$，故
$|T\cap Q|>(1-2^nt)\nu(Q)$。若 $a=0$，此式已排除小半径见证；若 $a>0$，取 $\gamma=1-2^nt$ 应用引理。周期 $T$ 的总测度无限，故先截取中心索引在 $[-N-1,N+1]^n$ 的 $T$ 副本；对 $[-N,N]^n$ 内半径小于 $1$ 的见证分子不变。应用引理、除以 $(2N)^n$ 并令 $N\to\infty$，边界因子趋于 $1$，得到每胞捕获量
\[
 b_{\rm small}\le \frac{4^nt}{1-2^nt}a
 =O(t^2u)=o(tu).\tag{P1}\label{eq:small-capture}
\]
这里用到的测度是 $D_\delta$ 上的内禀测度；没有将 Lebesgue 测度的 halo 包络直接移植到它上面。

若 $2\delta\le r\le1/2$，本孔完整包含于 $Q$，且无其它主孔进入。因此
$|D\cap Q|\ge m_0-a>t\ge t|Q|$，没有见证。对 $1/2\le r\le1$，令
\[
 s=\|x\|_\infty,\quad r_0=1-\delta-s,\quad
 A=\frac{m_0}{2\delta},\quad z_a=\frac{(3^n-1)a}{A}=O(u^2).
\]
每个边长至少 $1$ 的轴平行方体包含一个单位方体；周期性和 Fubini 表明任一单位方体内 $S$ 的体积恰为 $p$。故 $|S\cap Q|\ge p$。当 $r\le r_0$，唯一进入的原主孔完整包含，给出 $|D\cap Q|\ge q$。当 $r=r_0+z\in[r_0,1]$ 时，最近邻孔在一个坐标上的接触长度为 $z$，其余坐标至少覆盖本孔全宽，故 $|D_\delta\cap Q|\ge m_0+Az$。此方体至多交到 $3^n$ 个主孔副本；本孔删量至多 $a$，其它副本总删量至多 $(3^n-1)a$。于是
\[
 |D\cap Q|\ge q+A(z-z_a)_+.\tag{P2}\label{eq:neighbor-mass}
\]
该计数保留邻孔部分相交及删去部分恰在接触薄片上的最坏情形。

对充分小的 $t$，$q\le2m_0$、$\delta\le1/(6n)$。在 $z\ge z_a$ 上，函数
$F(z)=[q+A(z-z_a)]/[2(r_0+z)]^n$ 递增：导数分子除以 $A$ 后至少为
$r_0-(n-1)z+nz_a-nq/A\ge1-6n\delta\ge0$。对 $z\le z_a$，直接以分母 $[2(r_0+z_a)]^n$ 控制。连同 $r\le r_0$ 的情形可知
\[
 \frac{|D\cap Q(x,r)|}{|Q(x,r)|}\ge\frac q{[2(r_0+z_a)]^n}.
\]
故任何此范围内的见证都要求
\[
 \|x\|_\infty<Z:=1-\delta+z_a-
 \left(\frac q{2^nt}\right)^{1/n}.\tag{P3}\label{eq:Z-local}
\]

最后，对 $r\ge1$ 使用一维周期缺陷估计：若 $|x_j|\le\delta\le1/4$，则
\[
 \frac{|[x_j-r,x_j+r]\cap(\mathbb Z+[-\delta,\delta])|}{2r}
 \ge\frac32\delta.
\]
为核对这一维界，写 $r=k+\xi$，$k\ge1$ 为整数、$0\le\xi<1$，并记一维交长为 $\ell(x,r)$。周期性给 $\ell(x,r)=4k\delta+\ell(x,\xi)$。若 $\xi\le2\delta$，则 $\ell(x,\xi)\ge\xi$，所得比例至少为 $2\delta$；若 $\xi\ge2\delta$，则 $\ell(x,\xi)\ge2\delta$，比例至少为 $(2k+1)\delta/(k+1)\ge3\delta/2$。

逐坐标相乘给出 $|D_\delta\cap Q|/|Q|\ge(3\delta/2)^n$。交到 $Q$ 的删孔副本至多 $(2r+2)^n$ 个，故
$|T\cap Q|/|Q|\le(1+1/r)^na\le2^na$。因此
\[
 \frac{|D\cap Q|}{|Q|}\ge(3\delta/2)^n-2^na
 =((3/2)^n+o(1))t>t,
\]
大半径无见证。合并半径范围并允许两类捕获重叠，得到
\[
 G\le p+b_{\rm small}+\min\{m_0-a,[2Z_+]^n\},\qquad G\le q.
 \tag{P4}\label{eq:capture-total}
\]

现证系数。若归一化余项有一个不趋于 $-\infty$ 的子列，再取子列使其下界为常数。因 $n\ge2$，$t^2=o(tu)$；由 $\mathcal P_t(D)-1=G/(1-q)$ 及 $q\sim2^nt$ 得
\[
 G\ge2^nt-O(tu),\qquad q=m_0-a+p\ge2^nt-O(tu).
\]
写 $e=\delta/u-1=o(1)$，第二式推出 $e\ge-O(u)$。由 \eqref{eq:small-capture}、\eqref{eq:capture-total} 及 $p=O(tu)$，又有
$[2Z_+]^n\ge2^nt-O(tu)$，即 $Z\ge u-O(u^2)$。另一方面
\[
 \left(\frac q{2^nt}\right)^{1/n}
 =\frac\delta u\left(1+\frac{p-a}{m_0}\right)^{1/n}
 =1+e+O(u),
\]
故 \eqref{eq:Z-local} 给 $Z=-e+O(u)$，从而 $e\le O(u)$。于是 $\delta=u+O(u^2)$，并且 $Z\ge\delta-O(u^2)$。再代入 \eqref{eq:Z-local} 和 $z_a=O(u^2)$，得
\[
 \left(\frac q{2^nt}\right)^{1/n}\le1-2\delta+z_a+O(u^2)
 =1-2u+O(u^2),
\quad q\le2^nt-n2^{n+1}tu+O(tu^2).
\]
因此 $G/(1-q)\le q/(1-q)=q+O(t^2)$，除以 $tu$ 后误差趋于零。任一有固定下界的归一化余项子列满足定理上界；趋于 $-\infty$ 的子列也满足，定理得证。

为说明周期计算确实由有限测度输入逼近，固定 $t,D$，令
$E_N=E\cap[-N,N]^n$。其体积为 $(2N)^n(1-q)$。因 $1-t>1/2$，盒外任一中心方体至多一半体积落在盒内，故 $\{M_{c,n}\ind_{E_N}>1-t\}$ 位于盒内；单调性给出比值不超过 $\mathcal P_t(D)$。反向地，先只计周期 halo 中见证半径至多 $R$ 的点；离盒边界大于 $R$ 的点保留原见证，边界损失为 $O_{n,R}(N^{n-1})$。先令 $N\to\infty$，再令 $R\to\infty$，以半径截断水平集单调收敛得
\[
 \lim_{N\to\infty}\frac{|\{M_{c,n}\ind_{E_N}>1-t\}|}{|E_N|}
 =\mathcal P_t(D).
\]
所以结论排除这些有限测度截断序列在该局部周期族内提高二阶系数；不构成一般 $C_n$ 的全局上界。

\end{proof}
